\documentclass[11pt,letterpaper]{article}
\usepackage[margin=1in]{geometry}
\usepackage{amsmath}
\usepackage{listings}
\usepackage{xcolor}
\usepackage{hyperref}
\usepackage{graphicx}
\usepackage{booktabs}
\usepackage{enumitem}
\usepackage{parskip}

\definecolor{codegray}{rgb}{0.95,0.95,0.95}
\definecolor{commentgreen}{rgb}{0.13,0.54,0.13}
\definecolor{keywordblue}{rgb}{0.0,0.0,0.6}
\definecolor{stringred}{rgb}{0.6,0.0,0.0}

\lstset{
  language=Python,
  backgroundcolor=\color{codegray},
  basicstyle=\ttfamily\small,
  keywordstyle=\color{keywordblue}\bfseries,
  stringstyle=\color{stringred},
  commentstyle=\color{commentgreen}\itshape,
  breaklines=true,
  showstringspaces=false,
  frame=single,
  framesep=4pt,
  xleftmargin=6pt,
  xrightmargin=6pt,
  numbers=left,
  numberstyle=\tiny\color{gray},
  numbersep=6pt,
}

\title{\textbf{AutoResearch Loop: Implementation Report}\\
\large Feature 3 --- Autonomous Hyperparameter Search via Cyclic LangGraph}
\author{Matthew Torre \and Hayley Antczak \\
\small Stanford CS194W --- Team 26}
\date{May 2026}

\begin{document}
\maketitle

\tableofcontents
\newpage

% ─────────────────────────────────────────────────────────────────────────────
\section{Overview}

The AutoResearch Loop is the core iterative engine of the autonomous ML training
system. It accepts a \texttt{TrainingPlan} and \texttt{OrchestrationConfig}
from the Manager (Feature 0), and autonomously searches for hyperparameter
configurations that improve model performance---without any human intervention
between iterations.

The loop is implemented as a cyclic LangGraph \texttt{StateGraph} with nine
nodes and three conditional edges. Each iteration proposes one configuration
change, trains a model on Tinker's GPU infrastructure, evaluates the result,
and either commits or reverts the change. The process terminates when the budget
is exhausted, performance plateaus, or a hard iteration cap is reached. The
full experiment history is recorded in a JSONL research diary for
reproducibility and post-hoc analysis.

This report covers the architecture, design decisions, implementation details,
and the evaluation sub-feature (create/run/adapt eval suites).

% ─────────────────────────────────────────────────────────────────────────────
\section{Architecture}

\subsection{Graph Topology}

The loop is structured as a LangGraph \texttt{StateGraph(AutoResearchState)}.
The graph has a linear entry path and a cyclic core:

\begin{center}
\texttt{init} $\to$ \texttt{baseline} $\to$
$[$\texttt{propose} $\to$ \texttt{run} $\xrightarrow{\text{early-stop?}}$
\texttt{evaluate} $\xrightarrow{\text{keep/revert}}$
\texttt{log}$]^N$ $\to$ \texttt{END}
\end{center}

The graph is built once at startup by \texttt{build\_autoresearch\_graph()}
and executed via \texttt{graph.invoke(initial\_state)}. Table~\ref{tab:nodes}
summarises each node's responsibility.

\begin{table}[h]
\centering
\begin{tabular}{@{}lll@{}}
\toprule
\textbf{Node} & \textbf{Phase} & \textbf{Responsibility} \\
\midrule
\texttt{init}               & Setup    & Creates the eval suite \\
\texttt{baseline}           & Setup    & Runs unmodified script; establishes reference score \\
\texttt{propose}            & Cycle    & Asks Claude for one config change; patches \texttt{current.json} \\
\texttt{run}                & Cycle    & Submits Tinker job; blocks until done \\
\texttt{revert\_and\_continue} & Cycle & Early-stop path on catastrophic failure \\
\texttt{evaluate}           & Cycle    & Scores the model; computes delta vs.\ baseline \\
\texttt{keep}               & Cycle    & Commits patch; resets plateau counter \\
\texttt{revert}             & Cycle    & Rolls back patch; increments plateau counter \\
\texttt{log}                & Cycle    & Writes diary entry; adapts eval suite every 10 iters \\
\bottomrule
\end{tabular}
\caption{Node responsibilities in the AutoResearch graph.}
\label{tab:nodes}
\end{table}

\subsection{Conditional Edges}

Three conditional edge functions control the cycle:

\begin{itemize}[noitemsep]
  \item \textbf{\texttt{early\_stop\_edge}} (after \texttt{run}): returns
        \texttt{revert\_and\_continue} on catastrophic failure (NaN/inf loss,
        loss $> 10$, accuracy $< 0.01$); otherwise \texttt{evaluate}.
  \item \textbf{\texttt{decision\_edge}} (after \texttt{evaluate}): returns
        \texttt{keep} if \texttt{last\_delta.improved} is true, else
        \texttt{revert}.
  \item \textbf{\texttt{continue\_edge}} (after \texttt{log}): returns
        \texttt{\_\_end\_\_} when the budget is exhausted, the plateau counter
        reaches $\texttt{\_MAX\_NO\_IMPROVE}=3$, or total iterations reach
        $\texttt{\_MAX\_ITERATIONS}=20$; otherwise \texttt{propose}.
\end{itemize}

Each conditional edge is a pure function with no side effects, making them
independently testable with simple dict assertions.

\subsection{State Object}

All inter-node communication flows through a single \texttt{AutoResearchState}
TypedDict (defined in \texttt{src/types.py}). Nodes return only the fields
they update---never a full copy of state---which keeps LangGraph's
checkpointing overhead low and makes each node's contract explicit.

\subsection{File Layout}

\begin{center}
\begin{tabular}{@{}ll@{}}
\toprule
\textbf{File} & \textbf{Role} \\
\midrule
\texttt{src/autoresearch/autoresearch.py} & LangGraph graph, all nodes/edges, evaluator \\
\texttt{src/autoresearch/config.py}       & \texttt{TrainingConfig} dataclass + \texttt{BOUNDS} \\
\texttt{src/autoresearch/diff\_utils.py}  & Patch $\leftrightarrow$ diff-string serialisation \\
\texttt{src/autoresearch/proposer.py}     & Algorithmic proposal strategies \\
\texttt{src/autoresearch/loop.py}         & CLI scaffold for PROPOSE-only testing \\
\bottomrule
\end{tabular}
\end{center}

% ─────────────────────────────────────────────────────────────────────────────
\section{Design Decisions and Alternatives Considered}

\subsection{LangGraph vs.\ a Plain While Loop}

The most obvious alternative was a \texttt{while True} loop with
\texttt{if/else} branches for the keep/revert decision. This was prototyped
first. The problem: two structurally distinct revert paths (catastrophic
failure vs.\ no-improvement), the linear baseline-before-loop entry, budget
checking after every iteration, and eval suite adaptation every ten
iterations---all share mutable local state in a \texttt{while} loop, turning
control flow into nested conditionals that are hard to test or trace.

LangGraph solves this cleanly. The graph topology itself encodes the control
flow; reading \texttt{build\_autoresearch\_graph()} tells you the full cycle
without touching any node code. Nodes are pure functions testable in
isolation. LangGraph also provides checkpointing for free, which matters for
long-running GPU jobs that may need to resume after a preemption.

\subsection{Config-Level Patching vs.\ Script Patching}

The original spec described ``a single mutable training script plus patch
log,'' implying unified diffs of \texttt{train.py}. In practice, patching
Python source files with unified diffs is fragile: whitespace differences,
import ordering, and comment placement all cause patch application failures.
A corrupt \texttt{train.py} mid-run means the next job fails for reasons
unrelated to the hyperparameter being tested.

Patching \texttt{configs/current.json} instead avoids these issues entirely.
JSON is a structured format with a defined schema (\texttt{BOUNDS}), so every
patch can be validated before it touches disk. Writes are atomic via a
\texttt{.tmp} rename, so a crash mid-write never leaves a partial config.
The \texttt{.py} patching path is stubbed in \texttt{apply\_patch} with a
clear \texttt{NotImplementedError} for when it is eventually needed.

\subsection{Algorithmic Strategies vs.\ Claude-Only Proposals}

Two algorithmic proposal strategies (\texttt{RandomSearchProposalStrategy}
and \texttt{LocalPerturbationProposalStrategy}) were implemented alongside the
Claude API path. The strategies are used by the CLI scaffold for fast
testing without spending API budget. \texttt{LocalPerturbation} also
implements history-aware bias: it parses the diary's diff strings for
recently-reverted parameters and deprioritises them, which empirically
reduces wasted iterations in early testing.

Using Claude exclusively for proposals would have made the PROPOSE phase
impossible to test without an API key, and expensive to iterate on locally.
The dual-path design (strategies for testing, Claude for production) was
a deliberate trade-off.

\subsection{Stopping Conditions}

Three independent stopping conditions guard against runaway costs:

\begin{enumerate}[noitemsep]
  \item \textbf{Budget}: sum of \texttt{cost\_usd} across all diary entries
        $\geq$ \texttt{compute\_budget} from the \texttt{OrchestrationConfig}.
  \item \textbf{Convergence}: three consecutive non-improving iterations
        (\texttt{no\_improve\_streak} $\geq$ \texttt{\_MAX\_NO\_IMPROVE}).
  \item \textbf{Hard cap}: total iterations $\geq$ \texttt{\_MAX\_ITERATIONS}
        (20), regardless of improvement.
\end{enumerate}

Alternatives considered: Bayesian stopping criteria (too complex for the
evaluation timeline), fixed-iteration budgets (ignores budget and convergence),
and improvement-threshold stopping (too sensitive to noise in early iterations).

% ─────────────────────────────────────────────────────────────────────────────
\section{The Evaluator Sub-Feature}

\subsection{create\_eval\_suite}

Called by \texttt{init\_node}, this function determines which metrics matter
for the task type and whether LLM grading is needed. A \texttt{\_METRIC\_MAP}
maps task types to primary and supporting metrics (e.g.,
\texttt{text-classification} $\to$ F1 + accuracy + precision + recall). A
caller-specified \texttt{eval\_metric} in the \texttt{OrchestrationConfig}
overrides the default. LLM grading is enabled automatically for
high-complexity tasks.

\subsection{run\_evals}

This function was the main implementation gap at PR time. The full pipeline
requires loading a trained model from a Tinker artifact and running inference
on the test split---infrastructure that depends on Feature 4 (Tinker API)
being fully wired.

The implementation reads the \texttt{metrics.json} file saved by
\texttt{submit\_experiment} alongside each model checkpoint. This file
contains \texttt{train\_loss}, \texttt{val\_loss}, and \texttt{primary\_metric}
from the training run. These are mapped onto the eval suite's named metrics.
When no metrics file exists (e.g., the job failed before saving), a proxy
score of $\max(0, 1 - \text{val\_loss})$ is used.

For tasks with \texttt{use\_llm\_grading} enabled, up to five examples from
the test split are sampled and sent to Claude Haiku, which returns a quality
score and one-sentence critique. The final scalar is an 80/20 blend of the
training-derived score and the LLM grade:
\[
  s_{\text{final}} = 0.8 \cdot s_{\text{train}} + 0.2 \cdot s_{\text{llm}}
\]
This blending was chosen to keep the LLM grade as a signal rather than a
replacement: training metrics are more reliable for regression detection, but
LLM grading catches output quality issues that loss alone misses.

\subsection{adapt\_eval\_suite}

Called every ten iterations by \texttt{log\_node}. Collects the
\texttt{notes} fields from the last ten reverted iterations and asks Claude
Haiku to identify additional metrics or stress-test categories that would
surface those weaknesses. The returned metric names are merged into the eval
suite without duplicates. Once called, LLM grading is permanently enabled for
the suite, since the detected weaknesses imply heuristic metrics are
insufficient.

% ─────────────────────────────────────────────────────────────────────────────
\section{Key Implementation Details}

\subsection{Atomic Config Writes}

\texttt{apply\_patch} uses a write-then-rename pattern to prevent partial
writes from leaving a corrupt config on disk:

\begin{lstlisting}
tmp = path.with_suffix(".json.tmp")
tmp.write_text(json.dumps(config_dict, indent=2))
tmp.replace(path)
\end{lstlisting}

On POSIX systems, \texttt{Path.replace()} is guaranteed atomic. The original
file content is returned as a string so \texttt{revert\_patch} can restore it
byte-for-byte.

\subsection{Background Training Thread}

\texttt{submit\_experiment} launches a background \texttt{Thread} that runs
the Tinker SDK training loop. The thread stores its result in a module-level
dict keyed by job ID. \texttt{wait\_for\_experiment} joins the thread with a
configurable timeout (default 5 minutes) and raises \texttt{TimeoutError} if
the job does not complete in time.

\subsection{Early Stop Detection}

\texttt{check\_early\_stop} catches three failure modes before calling
\texttt{run\_evals}:

\begin{itemize}[noitemsep]
  \item \textbf{NaN/Inf loss}: \texttt{math.isnan} or \texttt{math.isinf}
        on \texttt{train\_loss} or \texttt{val\_loss}.
  \item \textbf{Exploding loss}: \texttt{train\_loss} $> 10.0$.
  \item \textbf{Accuracy collapse}: \texttt{primary\_metric} $< 0.01$.
\end{itemize}

Catching these before evaluation avoids spending API budget grading a clearly
failed model.

\subsection{Research Diary}

Every iteration appends one \texttt{IterationRecord} to
\texttt{outputs/logs/research\_diary.jsonl}. Each record contains the
iteration number, Claude's natural-language hypothesis, the git-style diff
of the config change, cost in USD, training metrics, and the
KEPT/REVERTED/PENDING decision. The diary serves as both an audit trail
and the context window that Claude reads when proposing the next change.

% ─────────────────────────────────────────────────────────────────────────────
\section{Testing}

Three test files cover the AutoResearch module:

\begin{itemize}
  \item \textbf{\texttt{test\_autoresearch.py}} --- Unit tests for every
        graph helper: \texttt{apply\_patch}/\texttt{revert\_patch} round-trips,
        \texttt{check\_early\_stop} on NaN/inf/exploding loss/accuracy collapse,
        \texttt{compare\_scores}, \texttt{decide\_keep\_or\_revert},
        \texttt{flag\_regression}, \texttt{log\_iteration}, all three
        \texttt{continue\_edge} stopping conditions, and the new
        \texttt{run\_evals} suite (seven tests covering missing/corrupt metrics
        files, the proxy fallback, per-metric dict completeness, and LLM
        grading skipping when the test split is empty).

  \item \textbf{\texttt{test\_propose\_infra.py}} --- Regression tests for
        eight specific bugs found during code review: NaN rejection in
        \texttt{apply\_patch}, \texttt{perturb\_value} handling
        \texttt{current=None}, seed increment across calls,
        \texttt{LocalPerturbation} reading regressed params from diff strings
        (not JSON), empty-patch guard in \texttt{run\_iteration}, corrupted
        diary line skipping, malformed diff line parsing, and atomic write
        leaving no temp file.

  \item \textbf{\texttt{test\_e2e\_propose.py}} --- End-to-end multi-iteration
        tests at three levels: mechanical (no API, fast), Claude API
        integration (skipped without \texttt{ANTHROPIC\_API\_KEY}), and a
        human-readable state trace.
\end{itemize}

All tests mock network calls with \texttt{unittest.mock.patch} and use
\texttt{tmp\_path} fixtures. The full suite runs in under one second:

\begin{center}
\texttt{python3 -m pytest tests/test\_autoresearch.py tests/test\_propose\_infra.py tests/test\_e2e\_propose.py -q}
\end{center}

% ─────────────────────────────────────────────────────────────────────────────
\section{Code Review Dimensions}

\subsection{Readability}
The graph topology in \texttt{build\_autoresearch\_graph()} serves as a
high-level map of the system. Each node is a short function (10--30 lines)
with a clear name. Conditional edges are named after the decision they
encode (\texttt{early\_stop\_edge}, \texttt{decision\_edge},
\texttt{continue\_edge}).

\subsection{Reusability}
\texttt{TrainingConfig} and \texttt{BOUNDS} are consumed by both the Decision
Engine and the AutoResearch Loop. \texttt{ProposalStrategy} is an abstract
base class, making it straightforward to plug in new strategies (e.g.,
Bayesian optimisation) without touching graph code.

\subsection{Security}
No user input reaches \texttt{subprocess} or \texttt{exec}. Config patches
are validated against \texttt{BOUNDS} before being written to disk.
The Claude API is called only with model-controlled system and user prompts
constructed entirely from typed Python objects.

\subsection{Performance}
Training jobs run in background threads, so the graph's main thread is
non-blocking. The research diary is appended incrementally (one JSONL line
per iteration) rather than rewritten, avoiding $O(n)$ writes per iteration.

\subsection{Robustness}
Three layers of failure handling: \texttt{check\_early\_stop} catches
catastrophic metrics before evaluation; \texttt{revert\_patch} restores the
config on any non-improvement; and \texttt{wait\_for\_experiment} raises
\texttt{TimeoutError} rather than hanging indefinitely. Corrupted diary lines
are silently skipped by \texttt{\_load\_history}.

% ─────────────────────────────────────────────────────────────────────────────
\section{Known Limitations and Future Work}

\begin{enumerate}
  \item \textbf{Inference-based evaluation}: \texttt{run\_evals} currently
        reads training metrics as a proxy. Full inference (loading a model
        from a Tinker artifact and running it over a held-out test split)
        requires Feature 4 (Tinker artifact download) to be fully wired. The
        structure is ready; it is a matter of plugging in the model-loading
        layer.

  \item \textbf{Cost Manager integration}: the \texttt{cost\_manager}
        parameter accepted by \texttt{invoke\_autoresearch\_graph} is not yet
        used; budget enforcement relies on summing diary costs. Wiring in the
        Cost Manager's polling API would provide real-time guardrails rather
        than post-hoc accounting.

  \item \textbf{Script-level patching}: \texttt{apply\_patch} raises
        \texttt{NotImplementedError} for \texttt{.py} files. Extending to
        unified diffs of \texttt{train.py} would allow the loop to explore
        architectural changes (e.g., adding gradient checkpointing), not just
        hyperparameter changes.

  \item \textbf{Thread safety}: \texttt{\_experiment\_results} and
        \texttt{\_experiment\_threads} are module-level dicts without locking.
        Sufficient for single-run use; would need a \texttt{threading.Lock}
        or a proper job queue for concurrent invocations.
\end{enumerate}

\end{document}
